\textbf{Introducción:} La deuda pública de Argentina exhibe episodios recurrentes de reestructuración e inestabilidad macroeconómica. Este trabajo evalúa la solvencia intertemporal de dicha deuda en el período 1996--2025 (30 años exactos) y simula escenarios estocásticos proyectados hasta 2035.

\textbf{Metodología:} Se aplica un protocolo econométrico secuencial sobre una ventana unificada de 120 observaciones trimestrales (1996T1--2025T4) construida mediante empalme histórico homogéneo, conservando el tramo 2004--2025 ($n=88$) como control de cobertura directa. Se estiman contrastes de raíz unitaria (ADF, KPSS, DF-GLS), quiebres estructurales endógenos (Zivot-Andrews, Bai-Perron) y cointegración de Johansen. La Función de Reacción Fiscal se modela mediante MCO Dinámicos (DOLS), Variables Instrumentales (IV-2SLS instrumentado con VIX y volatilidad bursátil del Bovespa), regresiones de umbral endógeno \citep{hansen1999,hansen2000} y un Modelo de Vectores con Corrección de Error (VECM). Como extensiones, se formula un SVAR restringido (Blanchard-Perotti/FMI), un análisis histórico de ultra-largo plazo del spread soberano (1983--2025) y un Análisis de Sostenibilidad de la Deuda (DSA) estocástico de Monte Carlo anclado en la matriz de covarianza del SVAR.

\textbf{Resultados:} La evidencia constata la ausencia de una regla fiscal lineal permanente: el coeficiente DOLS en muestra completa es positivo pero no significativo ($\rho=+0.0159$, $p=0.218$), y el VECM confirma la exogeneidad débil del superávit primario ($\alpha_{pb}=-0.00002$, $p=0.996$). En contraste, el modelo de umbral de Hansen en niveles identifica un régimen crítico en $\tau^*=2462.1$ puntos básicos ($p<0.001$), donde la reacción fiscal pasa de ser estadísticamente nula en normalidad ($\rho_1=-0.0044$, $p=0.665$) a fuertemente positiva bajo estrés financiero severo ($\rho_2=+0.0201$, $p=0.006$), revelando un patrón de disciplina fiscal gatillada por asfixia externa de crédito en lugar de fatiga fiscal. En el modelo IV-2SLS, el spread soberano ejerce una presión positiva y significativa sobre el ahorro primario ($\beta=0.0012$, $p=0.024$). La descomposición de varianza (FEVD) del SVAR demuestra que a 20 trimestres la inercia propia del stock acumulado explica el $54.21\%$ de la deuda, y el canal financiero-cambiario ($25.80\%$) supera al canal fiscal-real ($19.98\%$). La simulación estocástica DSA proyecta una probabilidad del $26.1\%$ de que la deuda supere el $100\%$ del PIB hacia 2035 (mediana de $82.8\%$).

\textbf{Discusión:} La solvencia de la deuda argentina no descansa en un compromiso continuo de ahorro presupuestario, sino en correcciones forzadas durante crisis de financiamiento y en canales de ajuste exógenos a la discreción fiscal (depreciación real, licuación inflacionaria y reestructuraciones soberanas), en consonancia con el paradigma de restricción externa. Las proyecciones DSA confirman que la trayectoria de mediano plazo está gobernada predominantemente por la volatilidad cambiaria y el costo de financiamiento externo antes que por el esfuerzo fiscal discrecional.

